\chapter{The Uninterpreted Opacity Zone}
\label{ch:opacity-zone}

Every type system draws a line: on one side lie the properties it can
decide; on the other, the properties it must approximate, defer, or
ignore.  In the $\PyObj$ theory the line has a precise geometric
characterization.  We call it the \emph{opacity boundary}.

This chapter gives the boundary a formal treatment.  We define three
interpretation levels, organise them into a lattice, prove that
soundness survives opacity while completeness degrades monotonically,
connect opacity to the \v{C}ech complex of Chapter~\ref{ch:cech}, and
derive the pipeline routing strategy that the implementation uses to
allocate proof effort efficiently.

%% ═══════════════════════════════════════════════════════════
\section{Interpretation Levels}
\label{sec:opacity-levels}
%% ═══════════════════════════════════════════════════════════

\begin{definition}[Opacity classification]
\label{def:opacity-levels}
Let $\Sigma$ be the signature of the $\PyObj$ Z3 theory.  Every
symbol $\sigma\in\Sigma$ belongs to exactly one of three
\emph{interpretation levels}:
\begin{enumerate}[label=(\roman*)]
  \item \textbf{Fully interpreted} ($\mathcal{I}_F$).
    The symbol is native to a decidable Z3 theory fragment
    (QF\_LIA, QF\_S, QF\_BV, QF\_BOOL).
    Z3 can reason about the \emph{value} of any closed ground
    term built from $\mathcal{I}_F$-symbols.

    \emph{Inhabitants:}
    $+, -, \times, \mathsf{//}, \mathsf{mod},
     <, \leq, >, \geq, =, \neq,
     \land, \lor, \lnot,
     \mathsf{Concat}, \mathsf{Length},
     \mathsf{LShift}, \mathsf{RShift},
     \mathsf{BitOr}, \mathsf{BitAnd}, \mathsf{BitXor}$.

  \item \textbf{Partially interpreted} ($\mathcal{I}_P$).
    The symbol is a Z3 \texttt{RecFunction} whose body references
    the $\PyObj$ ADT constructors.  Z3 can reason about
    \emph{structural} properties (constructor tags, Pair-chain
    shapes) but must unfold recursive definitions, making reasoning
    semi-decidable.

    \emph{Inhabitants:}
    $\mathsf{binop}, \mathsf{unop}, \mathsf{truthy},
     \mathsf{fold}, \mathsf{tag},
     \mathsf{GETITEM},
     \mathsf{mkint}, \mathsf{mkbool}, \mathsf{mkstr},
     \mathsf{mknone}, \mathsf{mkref}, \mathsf{mklist}$.

  \item \textbf{Uninterpreted} ($\mathcal{I}_U$).
    The symbol is a shared function symbol
    $\mathsf{py\_f} : \PyObj^k \to \PyObj$ with no built-in
    semantics.  Z3 knows only the congruence axiom:
    \[
      \vec{x} = \vec{y}
      \;\Longrightarrow\;
      \mathsf{py\_f}(\vec{x}) = \mathsf{py\_f}(\vec{y}).
    \]
    Optionally, equational axioms from
    $\mathcal{A}_{\mathsf{Py}}$ constrain the function further.

    \emph{Axiomatised inhabitants:}
    $\mathsf{sorted}, \mathsf{reversed}, \mathsf{list},
     \mathsf{tuple}, \mathsf{set}, \mathsf{frozenset},
     \mathsf{len}, \mathsf{sum},
     \mathsf{mut\_sort}, \mathsf{mut\_reverse}$.

    \emph{Bare (no axioms) inhabitants:}
    $\mathsf{dict}, \mathsf{any}, \mathsf{all},
     \mathsf{enumerate}, \mathsf{zip}, \mathsf{map},
     \mathsf{filter}, \mathsf{Counter}, \mathsf{defaultdict},
     \mathsf{deque}, \mathsf{type}, \mathsf{hash}, \mathsf{id},
     \mathsf{repr}, \mathsf{str}, \mathsf{iter}, \mathsf{next},
     \mathsf{print}, \mathsf{input}, \mathsf{open}$.
\end{enumerate}
\end{definition}

\begin{center}
\begin{tabular}{l l l l}
  \toprule
  \textbf{Level} & \textbf{Z3 power} & \textbf{Decidability}
    & \textbf{Example} \\
  \midrule
  Fully interpreted
    & Value reasoning
    & Decidable (QF\_LIA, QF\_S)
    & $3 + 4 = 7$ \\
  Partially interpreted
    & Structural reasoning
    & Semi-decidable (RecFunctions)
    & $\mathsf{truthy}(\mathsf{IntObj}(0)) = \mathsf{false}$ \\
  Uninterpreted
    & Congruence only
    & Decidable (QF\_UF)
    & $\mathsf{sorted}(x) = \mathsf{sorted}(x)$ \\
  \bottomrule
\end{tabular}
\end{center}


%% ═══════════════════════════════════════════════════════════
\section{The Opacity Lattice}
\label{sec:opacity-lattice}
%% ═══════════════════════════════════════════════════════════

\begin{definition}[Opacity ordering]
\label{def:opacity-order}
Define a partial order $\preceq$ on interpretation levels by
\[
  \mathcal{I}_F \;\preceq\; \mathcal{I}_P \;\preceq\; \mathcal{I}_U.
\]
The ordering reflects \emph{decreasing} Z3 reasoning power: a fully
interpreted symbol is the least opaque; an uninterpreted symbol is
the most.
\end{definition}

\begin{proposition}[Opacity lattice]
\label{prop:opacity-lattice}
The triple $(\{\mathcal{I}_F, \mathcal{I}_P, \mathcal{I}_U\},
\preceq)$ forms a bounded lattice with
\[
  \bot = \mathcal{I}_F, \qquad
  \top = \mathcal{I}_U, \qquad
  x \wedge y = \min(x,y), \qquad
  x \vee y = \max(x,y).
\]
\end{proposition}

\begin{definition}[Term opacity]
\label{def:term-opacity}
Let $t$ be a Z3 term over $\Sigma$.  The \emph{opacity} of $t$ is
\[
  \mathrm{opacity}(t) \;\coloneqq\;
  \bigvee_{\sigma \in \mathrm{symbols}(t)} \mathrm{level}(\sigma),
\]
the join (maximum) over the interpretation levels of every function
symbol appearing in $t$.
\end{definition}

\begin{corollary}
A term is fully interpreted if and only if every symbol in it belongs
to $\mathcal{I}_F$.  A single uninterpreted call suffices to raise
the whole term to $\mathcal{I}_U$.
\end{corollary}


%% ═══════════════════════════════════════════════════════════
\section{Formal Opacity Classification}
\label{sec:opacity-classification}
%% ═══════════════════════════════════════════════════════════

Given a Z3 term $t$, the \emph{opacity classifier} walks the term
tree and assigns each subterm a level.

\begin{definition}[Opacity profile]
\label{def:opacity-profile}
The \emph{opacity profile} of a term $t$ is the function
$\pi_t : \mathrm{Sub}(t) \to \{\mathcal{I}_F, \mathcal{I}_P,
\mathcal{I}_U\}$ defined recursively:
\begin{align}
  \pi_t(c) &= \mathcal{I}_F
    && \text{if } c \text{ is a constant or variable,} \\
  \pi_t(f(t_1,\ldots,t_k)) &= \mathrm{level}(f)
    \vee \bigvee_{i=1}^{k} \pi_t(t_i)
    && \text{otherwise.}
\end{align}
\end{definition}

The classifier in the implementation (see \texttt{OpacityClassifier}
in \texttt{proposals/opacity\_zone.py}) traverses the Z3 AST and
records:
\begin{enumerate}
  \item The per-node level $\pi_t(s)$ for every subterm $s$.
  \item The overall term opacity $\mathrm{opacity}(t)$.
  \item The set of \emph{blind spots}: maximal subterms at level
    $\mathcal{I}_U$ whose parent is at a lower level.
\end{enumerate}


%% ═══════════════════════════════════════════════════════════
\section{The Opacity Boundary}
\label{sec:opacity-boundary}
%% ═══════════════════════════════════════════════════════════

\begin{definition}[Opacity boundary]
\label{def:opacity-boundary}
Let $t$ be a term with $\mathrm{opacity}(t) = \mathcal{I}_U$.
An \emph{opacity boundary node} of $t$ is a subterm $s$ such that
$\pi_t(s) = \mathcal{I}_U$ but $s$ has a child $c$ with
$\pi_t(c) \prec \mathcal{I}_U$.  Equivalently, $s$ is the
shallowest point at which interpretation stops.
\end{definition}

\begin{theorem}[Opacity boundary theorem]
\label{thm:opacity-boundary}
Let $\mathsf{py\_f}$ be an uninterpreted function symbol.  Then:
\begin{enumerate}[label=(\alph*)]
  \item Z3 can prove
    $\mathsf{py\_f}(t_1) = \mathsf{py\_f}(t_2)$
    if and only if $t_1 = t_2$ is provable or an added axiom
    enables the deduction.
  \item Z3 \emph{cannot} prove $\mathsf{py\_f}(t) = v$ for any
    concrete value $v$ unless an axiom equates them.
  \item Adding the idempotence axiom
    $\forall x.\;\mathsf{py\_f}(\mathsf{py\_f}(x))
    = \mathsf{py\_f}(x)$ lifts $f$ partially out of the opacity
    zone: Z3 can simplify double applications but not compute
    output values.
\end{enumerate}
\end{theorem}

\begin{proof}
(a) follows directly from the congruence axiom scheme for
uninterpreted functions in the theory of equality with
uninterpreted functions (EUF).

(b) Because $\mathsf{py\_f}$ has no fixed interpretation, any model
that maps $\mathsf{py\_f}(t)$ to a value different from $v$ is
valid.  Hence the formula
$\mathsf{py\_f}(t) = v$ is satisfiable in one model and
falsifiable in another; Z3 returns $\mathsf{unknown}$ or
$\mathsf{sat}$ for the negation.

(c) The axiom
$\forall x.\;\mathsf{py\_f}(\mathsf{py\_f}(x)) = \mathsf{py\_f}(x)$
eliminates all models in which double application differs from
single application.  Combined with congruence, Z3 can now derive
$\mathsf{py\_f}(\mathsf{py\_f}(\mathsf{py\_f}(x))) =
\mathsf{py\_f}(x)$ and similar chains.  However, the axiom says
nothing about the output value for any single application, so (b)
still holds for non-double terms.
\end{proof}

\begin{remark}[Axioms as partially transparent windows]
\label{rem:axiom-windows}
Each semantic axiom ``opens a window'' in the opacity zone, allowing
Z3 to see one algebraic property of an otherwise opaque function.
The 52 axioms of $\mathcal{A}_{\mathsf{Py}}$ collectively open
enough windows for the benchmark suite while keeping most of each
function's behaviour opaque (and therefore decidable).
A window is \emph{algebraic} in the sense that it constrains the
function's behaviour on \emph{structured} inputs (compositions,
double applications) without revealing its behaviour on arbitrary
inputs.
\end{remark}


%% ═══════════════════════════════════════════════════════════
\section{Soundness under Opacity}
\label{sec:opacity-soundness}
%% ═══════════════════════════════════════════════════════════

\begin{theorem}[Soundness preservation]
\label{thm:opacity-soundness}
Let $\Gamma$ be the conjunction of all axioms in
$\mathcal{A}_{\mathsf{Py}}$ and let $\varphi$ be any formula over
$\Sigma$.  If $\Gamma \models \varphi$ (valid in all models of
$\Gamma$) then the corresponding Python semantic statement holds.
In particular, increasing the opacity of a symbol (removing axioms)
\emph{cannot introduce unsoundness}: it can only make Z3 fail to
prove a valid statement.
\end{theorem}

\begin{proof}
Soundness of the $\PyObj$ theory was established in
Theorem~\ref{thm:z3-soundness}.  Each axiom in
$\mathcal{A}_{\mathsf{Py}}$ is a true statement about Python
semantics (verified by ground testing in \S\ref{sec:recfn-tests}).
Removing axioms weakens $\Gamma$, so
$\Gamma' \subseteq \Gamma$ implies
$\{\varphi : \Gamma' \models \varphi\}
  \subseteq
\{\varphi : \Gamma \models \varphi\}$.
No new theorems are introduced; some existing ones become
unprovable.
\end{proof}

\begin{corollary}
Any equivalence that Z3 proves under the current axiom set is a
true Python equivalence, regardless of how many symbols are
uninterpreted.
\end{corollary}


%% ═══════════════════════════════════════════════════════════
\section{The Completeness Gap Theorem}
\label{sec:completeness-gap}
%% ═══════════════════════════════════════════════════════════

\begin{definition}[Provable equivalence set]
\label{def:provable-equiv}
For a set of axioms $\Gamma$ over signature $\Sigma$, define
\[
  \mathrm{Eq}(\Gamma)
  \;\coloneqq\;
  \bigl\{ (t_1, t_2) \in \mathrm{Term}(\Sigma)^2
    : \Gamma \models t_1 = t_2 \bigr\}.
\]
\end{definition}

\begin{theorem}[Completeness gap]
\label{thm:completeness-gap}
Let $\Gamma_0 \subseteq \Gamma_1 \subseteq \cdots \subseteq
\Gamma_n$ be an increasing chain of axiom sets (each $\Gamma_{i+1}$
adds axioms for one more builtin).  Then:
\[
  \mathrm{Eq}(\Gamma_0)
  \;\subseteq\;
  \mathrm{Eq}(\Gamma_1)
  \;\subseteq\;
  \cdots
  \;\subseteq\;
  \mathrm{Eq}(\Gamma_n)
  \;\subseteq\;
  \mathrm{Eq}_{\mathsf{Py}},
\]
where $\mathrm{Eq}_{\mathsf{Py}}$ is the set of all true Python
equivalences.  Each inclusion is strict in general, and the final
gap $\mathrm{Eq}(\Gamma_n) \subsetneq \mathrm{Eq}_{\mathsf{Py}}$
is unavoidable (by Rice's theorem).
\end{theorem}

\begin{proof}
Monotonicity: adding axioms can only increase the set of
provable consequences, so each inclusion holds.

Strictness of intermediate inclusions: consider two programs
$f(x) = \mathsf{sorted}(\mathsf{list}(x))$ and $g(x) = \mathsf{sorted}(x)$.
With $\Gamma_0$ (no axioms for \textsf{sorted} or \textsf{list}),
Z3 cannot prove $f = g$.  With $\Gamma_1$ containing the axiom
$\forall x.\;\mathsf{sorted}(\mathsf{list}(x)) = \mathsf{sorted}(x)$,
Z3 can.  Hence $\mathrm{Eq}(\Gamma_0) \subsetneq \mathrm{Eq}(\Gamma_1)$.

Unavoidability of the final gap: by Rice's theorem, the set of
semantically equivalent program pairs is undecidable, so no finite
axiom set can capture all of $\mathrm{Eq}_{\mathsf{Py}}$.
\end{proof}

\begin{proposition}[Monotone opacity--completeness trade-off]
\label{prop:opacity-completeness}
Define the \emph{opacity score} of a pair $(t_1, t_2)$ as
$\max(\mathrm{opacity}(t_1), \mathrm{opacity}(t_2))$.
As the opacity score increases from $\mathcal{I}_F$ to
$\mathcal{I}_U$, the conditional probability that Z3 can prove
$t_1 = t_2$ (given that $t_1$ and $t_2$ are truly equivalent)
decreases monotonically.
\end{proposition}


%% ═══════════════════════════════════════════════════════════
\section{Opacity and the \v{C}ech Complex}
\label{sec:opacity-cech}
%% ═══════════════════════════════════════════════════════════

Recall from Chapter~\ref{ch:cech} that the \v{C}ech complex of two
programs $f, g$ is built from \emph{fibers} over input-space
patches, and equivalence is witnessed by the vanishing of
$H^1(\mathcal{U}; \mathcal{F}_{f-g})$.

\begin{definition}[Blind spot]
\label{def:blind-spot}
A \emph{blind spot} of the \v{C}ech equivalence checker is a fiber
$(U_i, t_f, t_g)$ where both $t_f$ and $t_g$ have opacity
$\mathcal{I}_U$ and share no common substructure that Z3 can
exploit.  In a blind spot, Z3 returns $\mathsf{unknown}$ or
(worse) $\mathsf{unsat}$ vacuously because it is free to choose
\emph{any} interpretation for the uninterpreted functions.
\end{definition}

\begin{theorem}[Vacuous equivalence]
\label{thm:vacuous-equiv}
If $t_f = \mathsf{py\_h}(s)$ and $t_g = \mathsf{py\_k}(s')$ for
distinct uninterpreted symbols $\mathsf{py\_h}, \mathsf{py\_k}$
with no relating axiom, then Z3 can find a model where
$\mathsf{py\_h} = \mathsf{py\_k}$ (mapping them to the same
function), yielding a spurious $\mathsf{unsat}$ for the
disequality $t_f \neq t_g$.
This is \emph{vacuously true} equivalence and must be rejected.
\end{theorem}

\begin{proof}
In the theory of uninterpreted functions, function symbols are
universally quantified over all possible interpretations.
Nothing prevents a model from interpreting two distinct symbols
identically.  Hence the solver can always satisfy $t_f = t_g$
by collapsing the two symbols, even when the Python functions
compute different things.
\end{proof}

The implementation guards against this via the
\texttt{terms\_same\_opacity} check (see~\S\ref{sec:pipeline}):
an $\mathsf{unsat}$ result is trusted only when both $t_f$ and
$t_g$ are fully interpreted (built exclusively from RecFunctions
with known semantics like $\mathsf{binop}$, $\mathsf{truthy}$,
$\mathsf{fold}$).

\begin{corollary}[Opacity creates \v{C}ech blind spots]
\label{cor:opacity-blind}
The number of blind spots in the \v{C}ech complex is bounded below by
the number of fibers $(U_i, t_f, t_g)$ where
$\mathrm{opacity}(t_f) \vee \mathrm{opacity}(t_g) = \mathcal{I}_U$
and both terms reference distinct uninterpreted symbols.
\end{corollary}


%% ═══════════════════════════════════════════════════════════
\section{Opacity and Path Search}
\label{sec:opacity-path}
%% ═══════════════════════════════════════════════════════════

When Z3 cannot determine fiber equivalence (blind spots), the
pipeline falls back to \emph{path search}: an explicit search
through the equational theory $\mathcal{A}_{\mathsf{Py}}$ for
a chain of rewrites connecting $t_f$ to $t_g$.

\begin{definition}[Opacity-triggered path search]
\label{def:opacity-path}
Let $(U_i, t_f, t_g)$ be a fiber with
$\mathrm{opacity}(t_f) \vee \mathrm{opacity}(t_g) = \mathcal{I}_U$.
The path search procedure:
\begin{enumerate}
  \item Identifies every axiom in $\mathcal{A}_{\mathsf{Py}}$ that
    mentions a symbol appearing in $t_f$ or $t_g$.
  \item Constructs the rewrite graph $\mathcal{R}$ with terms as
    nodes and axiom applications as directed edges.
  \item Searches $\mathcal{R}$ (BFS or iterative deepening) for a
    path from $t_f$ to $t_g$.
  \item If a path is found, the sequence of axiom names constitutes
    a \emph{proof term} for $t_f = t_g$.
\end{enumerate}
\end{definition}

\begin{proposition}[Complementarity]
\label{prop:complementarity}
Z3 and path search are complementary along the opacity axis:
\begin{itemize}
  \item \emph{Low opacity} ($\mathcal{I}_F$): Z3 decides the fiber
    in milliseconds; path search is unnecessary.
  \item \emph{Medium opacity} ($\mathcal{I}_P$): Z3 may need
    RecFunction unfolding, which can time out; path search provides
    a backup by exploiting algebraic structure.
  \item \emph{High opacity} ($\mathcal{I}_U$): Z3 can only apply
    congruence; path search through axioms is the primary proof
    method.
\end{itemize}
\end{proposition}


%% ═══════════════════════════════════════════════════════════
\section{The Opacity-Aware Pipeline}
\label{sec:pipeline}
%% ═══════════════════════════════════════════════════════════

The implementation routes each program pair through the verification
pipeline based on an opacity pre-analysis.

\begin{definition}[Pipeline routing]
\label{def:pipeline-routing}
Given compiled terms $t_f, t_g$ for programs $f, g$:
\begin{enumerate}
  \item Compute $\alpha \coloneqq \mathrm{opacity}(t_f) \vee
    \mathrm{opacity}(t_g)$.
  \item Route:
    \begin{itemize}
      \item $\alpha = \mathcal{I}_F$: \textbf{Z3 direct}.  Submit
        $t_f \neq t_g$ to Z3 with a short timeout.
        The result is definitive.
      \item $\alpha = \mathcal{I}_P$: \textbf{Z3 + path search}.
        Try Z3 first; on timeout, run path search.
        Trust Z3 $\mathsf{unsat}$ results (both terms are
        structurally interpreted).
      \item $\alpha = \mathcal{I}_U$: \textbf{Path search +
        spec match}.  Skip Z3 for equivalence (vacuous risk);
        use Z3 only for \emph{counterexample} search.
        Rely on axiom-based path search and specification matching.
    \end{itemize}
\end{enumerate}
\end{definition}

\begin{lstlisting}[caption={Opacity-aware routing (pseudocode)}]
def route_pair(tf, tg, axioms):
    alpha = max(opacity(tf), opacity(tg))
    if alpha == FULLY_INTERPRETED:
        return z3_decide(tf, tg)
    elif alpha == PARTIALLY_INTERPRETED:
        r = z3_decide(tf, tg, timeout=5000)
        if r is not None:
            return r
        return path_search(tf, tg, axioms)
    else:  # UNINTERPRETED
        cex = z3_find_counterexample(tf, tg)
        if cex is not None:
            return INEQUIVALENT
        return path_search_or_spec_match(tf, tg, axioms)
\end{lstlisting}

\begin{remark}[The opacity check in the \v{C}ech verifier]
\label{rem:opacity-check}
When Z3 reports $\mathsf{unsat}$ for the disequality
$t_f \neq t_g$, the proof may be vacuously true if either term
involves uninterpreted functions.  Z3 is free to choose \emph{any}
interpretation for uninterpreted function symbols, so it can
make them equal by fiat.  The \texttt{terms\_same\_opacity}
function checks that both terms are \emph{fully interpreted}
(built only from RecFunctions with known semantics such as
$\mathsf{binop}$, $\mathsf{truthy}$, $\mathsf{fold}$).
Only then is the $\mathsf{unsat}$ result trusted.

When the opacity check fails, the fiber is marked
\emph{inconclusive} rather than equivalent, and the pipeline
defers to path search.
\end{remark}


%% ═══════════════════════════════════════════════════════════
\section{Opacity Reduction Strategies}
\label{sec:opacity-reduction}
%% ═══════════════════════════════════════════════════════════

Given a term $t$ with high opacity, several strategies can
\emph{reduce} the opacity, bringing more of $t$'s behaviour into
Z3's reasoning scope.

\subsection{Axiomatisation}

\begin{definition}[Axiom-based reduction]
\label{def:axiom-reduction}
Adding an equational axiom for an uninterpreted symbol $\mathsf{py\_f}$
does not change $\mathsf{py\_f}$'s classification (it remains in
$\mathcal{I}_U$) but enlarges $\mathrm{Eq}(\Gamma)$ by enabling
new derivations.  Each axiom opens a ``window'' through which Z3
can see one algebraic property.

\emph{The fundamental trade-off}:  More axioms $\Longrightarrow$
larger $\mathrm{Eq}(\Gamma)$ (more completeness) but more
universally quantified formulas for Z3 to instantiate
(potential timeouts).  In practice, the 52 axioms of
$\mathcal{A}_{\mathsf{Py}}$ hit a sweet spot: enough to handle
the benchmark suite without causing Z3 timeouts.
\end{definition}

The axioms from \texttt{theory.py}'s \texttt{semantic\_axioms()}
method include:
\begin{itemize}
  \item \textbf{S4}: $\mathsf{sorted}$ idempotent —
    $\forall x.\;\mathsf{sorted}(\mathsf{sorted}(x)) = \mathsf{sorted}(x)$.
  \item \textbf{S7}: $\mathsf{reversed}$ involution —
    $\forall x.\;\mathsf{reversed}(\mathsf{reversed}(x)) = x$.
  \item \textbf{I3--I6}: Collection idempotence —
    $\forall x.\; c(c(x)) = c(x)$ for $c \in
    \{\mathsf{list}, \mathsf{set}, \mathsf{frozenset}, \mathsf{tuple}\}$.
  \item \textbf{S6, S8}: $\mathsf{len}$ preservation —
    $\forall x.\;\mathsf{len}(d(x)) = \mathsf{len}(x)$ for
    $d \in \{\mathsf{sorted}, \mathsf{reversed}, \mathsf{list},
    \mathsf{tuple}\}$.
  \item \textbf{Sum invariance}: $\mathsf{sum}$ is
    permutation-invariant —
    $\forall x.\;\mathsf{sum}(d(x)) = \mathsf{sum}(x)$ for
    $d \in \{\mathsf{sorted}, \mathsf{reversed}, \mathsf{list}\}$.
  \item \textbf{S9}: $\mathsf{sorted}$ absorbs $\mathsf{reversed}$ —
    $\forall x.\;\mathsf{sorted}(\mathsf{reversed}(x))
    = \mathsf{sorted}(x)$.
  \item \textbf{Absorption}: $\mathsf{sorted}$ absorbs
    $\mathsf{list}/\mathsf{tuple}$; $\mathsf{set}$ absorbs
    $\mathsf{sorted}/\mathsf{reversed}/\mathsf{list}/\mathsf{tuple}$.
  \item \textbf{Mutation bridges}: $\mathsf{mut\_sort}(\mathsf{list}(x))
    = \mathsf{sorted}(x)$;
    $\mathsf{mut\_reverse}(\mathsf{list}(x))
    = \mathsf{list}(\mathsf{reversed}(x))$.
\end{itemize}

\subsection{Ground Evaluation}

\begin{definition}[Ground evaluation]
\label{def:ground-eval}
If a subterm $s$ of $t$ is \emph{closed} (contains no free
variables), evaluate $s$ by executing the corresponding Python
expression on concrete inputs and replacing $s$ with a fully
interpreted constant.

Formally, let $\mathrm{eval} : \mathrm{ClosedTerm}(\Sigma) \to
\PyObj$ be the Python evaluation function.  Replace every maximal
closed subterm $s$ with
$\iota(\mathrm{eval}(s))$, where $\iota$ is the embedding into
$\PyObj$.  The resulting term has strictly lower opacity.
\end{definition}

\begin{example}
The term $\mathsf{py\_sorted}(\mathsf{mklist}(3, 1, 2))$ is closed.
Ground evaluation yields $\mathsf{mklist}(1, 2, 3)$, replacing
an uninterpreted application with a fully interpreted Pair-chain.
Opacity drops from $\mathcal{I}_U$ to $\mathcal{I}_P$.
\end{example}

\subsection{Specification Matching}

\begin{definition}[Specification identification]
\label{def:spec-match}
A \emph{specification pattern} is a pair $(\pi, t_{\mathrm{canon}})$
where $\pi$ is a syntactic pattern over $\Sigma$ (e.g.,
``$\mathsf{sorted}(\mathsf{list}(\mathsf{set}(\Box)))$'') and
$t_{\mathrm{canon}}$ is a canonical equivalent term.  If a term $t$
matches $\pi$ with substitution $\theta$, replace $t$ with
$t_{\mathrm{canon}}\theta$.  This bypasses Z3 entirely via
pre-verified algebraic identities.
\end{definition}

\subsection{Normalisation}

\begin{definition}[Opacity-reducing normalisation]
\label{def:normalise}
Apply rewrite rules that eliminate uninterpreted subterms:
\begin{align}
  \mathsf{sorted}(\mathsf{list}(x)) &\;\longrightarrow\;
    \mathsf{sorted}(x), \\
  \mathsf{reversed}(\mathsf{reversed}(x)) &\;\longrightarrow\; x, \\
  \mathsf{list}(\mathsf{list}(x)) &\;\longrightarrow\;
    \mathsf{list}(x), \\
  \mathsf{set}(\mathsf{sorted}(x)) &\;\longrightarrow\;
    \mathsf{set}(x).
\end{align}
Each rule is oriented so that the right-hand side has equal or
lower opacity depth (the longest chain of nested uninterpreted
calls).
\end{definition}


%% ═══════════════════════════════════════════════════════════
\section{The Fundamental Trade-Off}
\label{sec:tradeoff}
%% ═══════════════════════════════════════════════════════════

\begin{theorem}[Axiom complexity trade-off]
\label{thm:tradeoff}
Let $|\Gamma|$ denote the number of universally quantified axioms
in the Z3 context.  Then:
\begin{enumerate}[label=(\alph*)]
  \item $|\mathrm{Eq}(\Gamma)|$ is non-decreasing in $|\Gamma|$
    (more axioms, more provable equivalences).
  \item The expected Z3 solving time is non-decreasing in $|\Gamma|$
    (more axioms, more quantifier instantiations).
  \item There exists a threshold $|\Gamma^*|$ beyond which Z3
    routinely times out, making the additional axioms useless in
    practice.
\end{enumerate}
\end{theorem}

\begin{proof}
(a) is immediate from monotonicity of logical consequence.
(b) follows from the MBQI (model-based quantifier instantiation)
algorithm: each universally quantified axiom generates candidate
instantiations, and the search space grows combinatorially.
(c) is empirical: in the benchmark suite, adding more than
$\sim\!60$ universally quantified axioms causes Z3 to exceed its
5-second timeout on $>20\%$ of fibers.
\end{proof}

\begin{remark}[The 52-axiom sweet spot]
The current $\mathcal{A}_{\mathsf{Py}}$ contains 52 axioms.
This was determined empirically: enough to handle the
sorted/reversed/set/list/tuple/len/sum ecosystem while staying
well within Z3's timeout budget.  Adding axioms for
$\mathsf{dict}$, $\mathsf{Counter}$, or $\mathsf{map}$ would
increase coverage but risks pushing Z3 past the timeout cliff.
\end{remark}


%% ═══════════════════════════════════════════════════════════
\section{Worked Examples}
\label{sec:opacity-examples}
%% ═══════════════════════════════════════════════════════════

\subsection{Fully Opaque: \texttt{sorted(list(set(x)))}}

Consider comparing \verb|sorted(list(set(x)))| with \verb|sorted(x)|.
The compiler produces:

\begin{lstlisting}[caption={Fully opaque chain}]
p0 = Const('p0', S)          # input x
t1 = py_set(p0)              # set(x)
t2 = py_list(t1)             # list(set(x))
t3 = py_sorted(t2)           # sorted(list(set(x)))

t4 = py_sorted(p0)           # sorted(x)
\end{lstlisting}

\textbf{Opacity analysis.}  Every function symbol in $t_3$ is in
$\mathcal{I}_U$.  The opacity profile:
\[
  \pi(p_0) = \mathcal{I}_F, \quad
  \pi(\mathsf{set}(p_0)) = \mathcal{I}_U, \quad
  \pi(\mathsf{list}(\cdot)) = \mathcal{I}_U, \quad
  \pi(\mathsf{sorted}(\cdot)) = \mathcal{I}_U.
\]
The opacity boundary is at $\mathsf{set}(p_0)$: the innermost
uninterpreted call.

\textbf{Without axioms.}  Z3 cannot prove $t_3 = t_4$ because
$\mathsf{sorted}(\mathsf{list}(\mathsf{set}(p_0)))$ and
$\mathsf{sorted}(p_0)$ are syntactically distinct.  No congruence
chain connects them.

\textbf{With axioms.}  Adding the absorption axioms:
\begin{align}
  \forall x.\; \mathsf{sorted}(\mathsf{list}(x))
    &= \mathsf{sorted}(x)
    \tag{sorted absorbs list} \\
  \forall x.\; \mathsf{set}(\mathsf{set}(x))
    &= \mathsf{set}(x)
    \tag{I3: set idempotent}
\end{align}
The chain reduces:
$\mathsf{sorted}(\mathsf{list}(\mathsf{set}(p_0)))
\xrightarrow{\text{absorb list}}
\mathsf{sorted}(\mathsf{set}(p_0))$.

This is \emph{not} equal to $\mathsf{sorted}(p_0)$ because
$\mathsf{set}$ removes duplicates (changes the multiset).
Z3 correctly reports \textbf{not equivalent}.

\subsection{Partially Opaque: \texttt{sum(range(n))} vs \texttt{n*(n-1)//2}}

\begin{lstlisting}[caption={Mixed opacity levels}]
# Program f: sum(range(n))
p0 = Const('p0', S)          # input n (IntObj)
t1 = py_range(p0)            # range(n): uninterpreted
t2 = py_sum(t1)              # sum(range(n)): uninterpreted

# Program g: n*(n-1)//2
t3 = binop(MUL, p0,          # n * (n-1): fully interpreted
       binop(SUB, p0, mkint(1)))
t4 = binop(FLOORDIV, t3,     # ... // 2: fully interpreted
       mkint(2))
\end{lstlisting}

\textbf{Opacity analysis.}
$\mathrm{opacity}(t_2) = \mathcal{I}_U$
(both \textsf{range} and \textsf{sum} are uninterpreted).
$\mathrm{opacity}(t_4) = \mathcal{I}_F$
(only arithmetic).
The pair opacity is $\mathcal{I}_U$.

\textbf{Pipeline routing.}
The pair is routed to path search + spec match.
Z3 is used only for counterexample search (which fails because
the programs are equivalent on non-negative integers).
Path search also fails (no axioms connect $\mathsf{sum}$ and
$\mathsf{range}$ to arithmetic).  The pair is reported as
\emph{inconclusive} --- correctly, because a proof requires
mathematical induction beyond the equational theory.

\subsection{Fully Interpreted: \texttt{x + 0} vs \texttt{x}}

\begin{lstlisting}[caption={Fully interpreted}]
# Program f: x + 0
t1 = binop(ADD, p0, mkint(0))

# Program g: x (identity)
t2 = p0
\end{lstlisting}

\textbf{Opacity analysis.}
$\mathrm{opacity}(t_1) = \mathcal{I}_F$ (all symbols in QF\_LIA).
$\mathrm{opacity}(t_2) = \mathcal{I}_F$ (a variable).
Pair opacity is $\mathcal{I}_F$.

\textbf{Pipeline routing.}
Routed to Z3 direct.  The solver expands
$\mathsf{binop}(\mathsf{ADD}, \mathsf{IntObj}(n),
\mathsf{IntObj}(0))$ to $\mathsf{IntObj}(n + 0) = \mathsf{IntObj}(n)$,
which QF\_LIA decides immediately.
Result: \textbf{equivalent}.

Note that the solver must unfold the $\mathsf{binop}$ RecFunction,
which involves the $\PyObj$ ADT (partially interpreted), but
\emph{after} unfolding the guard $\mathsf{is\_IntObj}(a) \land
\mathsf{is\_IntObj}(b)$ holds and the body reduces to pure
integer arithmetic.  This is why the opacity of a $\mathsf{binop}$
application is context-dependent: on integer arguments it
effectively reduces to $\mathcal{I}_F$.


%% ═══════════════════════════════════════════════════════════
\section{Implementation Mapping}
\label{sec:opacity-impl}
%% ═══════════════════════════════════════════════════════════

The opacity classification maps directly to the code in
\texttt{theory.py} and \texttt{compiler.py}.

\subsection{Fully Interpreted Operations in \texttt{theory.py}}

The \texttt{\_define\_binop} method encodes integer arithmetic,
string concatenation, boolean logic, and comparisons directly
in Z3's native sorts:

\begin{lstlisting}[caption={Fully interpreted: arithmetic (theory.py)}]
int_r = If(op == ADD, S.IntObj(ai + bi),
        If(op == SUB, S.IntObj(ai - bi),
        If(op == MUL, S.IntObj(ai * bi),
        ...)))
\end{lstlisting}

These terms use \texttt{IntSort()} arithmetic ($+, -, \times$),
\texttt{BoolSort()} comparisons ($<, \leq, =, \neq$), and
\texttt{StringSort()} operations (\texttt{Concat}, \texttt{Length}).
Z3 has complete decision procedures for all of them.

\subsection{Partially Interpreted: RecFunctions}

The five RecFunctions (\texttt{binop}, \texttt{unop}, \texttt{truthy},
\texttt{fold}, \texttt{tag}) are structurally interpreted: their
bodies reference the $\PyObj$ ADT constructors and dispatch on
constructor tags.  Z3 can unfold them but doing so is
semi-decidable (the solver may not terminate on recursive inputs).

\subsection{Uninterpreted: Shared Function Symbols}

The \texttt{shared\_fn} method in \texttt{theory.py} creates
uninterpreted function symbols:

\begin{lstlisting}[caption={Uninterpreted shared symbols (theory.py)}]
def shared_fn(self, name, arity):
    key = (name, arity)
    if key not in self._shared_fns:
        sorts = [self.S] * arity + [self.S]
        self._shared_fns[key] = Function(f'py_{name}', *sorts)
    return self._shared_fns[key]
\end{lstlisting}

The \texttt{compiler.py} module determines which builtins produce
shared symbols vs native Z3 terms.  The \texttt{SHARED\_BUILTINS}
set lists all builtins compiled to uninterpreted symbols:

\begin{lstlisting}[caption={Shared builtins (compiler.py)}]
SHARED_BUILTINS = frozenset({
    'sorted', 'reversed', 'list', 'tuple', 'set', 'frozenset',
    'dict', 'sum', 'any', 'all', 'enumerate', 'zip', 'map',
    'filter', 'type', 'hash', 'id', 'repr', 'str', 'iter',
    'next', 'print', 'input', 'open', 'ord', 'chr', ...
})
\end{lstlisting}

By contrast, calls to \texttt{abs}, \texttt{int}, \texttt{bool},
\texttt{len}, \texttt{max(a,b)}, \texttt{min(a,b)} are compiled
to native Z3 operations:

\begin{lstlisting}[caption={Native compilation (compiler.py)}]
if name == 'abs' and n == 1: return T.unop(ABS_, args[0])
if name == 'int' and n == 1: return T.unop(INT_, args[0])
if name == 'len' and n == 1: return T.unop(LEN_, args[0])
if name == 'max' and n == 2: return T.binop(MAX, args[0], args[1])
\end{lstlisting}

These bypass the shared-symbol mechanism entirely and produce
partially- or fully-interpreted terms.


%% ═══════════════════════════════════════════════════════════
\section{Opacity Metrics}
\label{sec:opacity-metrics}
%% ═══════════════════════════════════════════════════════════

\begin{definition}[Opacity score]
\label{def:opacity-score}
The \emph{opacity score} of a term $t$ is
\[
  \mathrm{score}(t) \;\coloneqq\;
  \frac{|\{s \in \mathrm{Sub}(t) : \pi_t(s) = \mathcal{I}_U\}|}
       {|\mathrm{Sub}(t)|},
\]
the fraction of subterms that are uninterpreted.
\end{definition}

\begin{definition}[Opacity depth]
\label{def:opacity-depth}
The \emph{opacity depth} of a term $t$ is the maximum nesting
depth of uninterpreted function applications:
\[
  \mathrm{depth}_U(t) \;\coloneqq\;
  \begin{cases}
    0 & \text{if } t \text{ is a variable or constant,} \\
    1 + \max_i \mathrm{depth}_U(t_i)
      & \text{if } t = f(t_1,\ldots,t_k)
        \text{ and } f \in \mathcal{I}_U, \\
    \max_i \mathrm{depth}_U(t_i)
      & \text{otherwise.}
  \end{cases}
\]
\end{definition}

\begin{definition}[Maximum opaque chain length]
\label{def:opaque-chain}
The \emph{maximum opaque chain length} of $t$ is the length of the
longest path in $t$'s syntax tree that passes only through
$\mathcal{I}_U$-level function applications.  This measures how
many uninterpreted functions are composed without any interpreted
``checkpoint'' where Z3 can reason about intermediate values.
\end{definition}

\begin{proposition}[Completeness estimate]
\label{prop:completeness-estimate}
Let $\alpha = \mathrm{score}(t_f) \vee \mathrm{score}(t_g)$.
Define the \emph{estimated Z3 success probability}
\[
  P_{\mathrm{Z3}}(\alpha) \;\coloneqq\; (1 - \alpha)^2.
\]
This is a rough heuristic: when $\alpha = 0$ (fully interpreted),
$P = 1$; when $\alpha = 1$ (fully opaque), $P = 0$.  The quadratic
shape reflects the fact that \emph{both} terms must be interpreted
for Z3 to give a reliable answer.
\end{proposition}


%% ═══════════════════════════════════════════════════════════
\section{Summary}
\label{sec:opacity-summary}
%% ═══════════════════════════════════════════════════════════

The opacity zone is not a defect but a design feature.  By keeping
most of Python's standard library uninterpreted, the $\PyObj$
theory remains in decidable fragments (QF\_UF + QF\_LIA) while
the equational axioms $\mathcal{A}_{\mathsf{Py}}$ selectively open
windows into the opaque functions.  The opacity-aware pipeline
routes each pair to the appropriate proof strategy, and the
\texttt{terms\_same\_opacity} guard prevents vacuous equivalence
proofs.

The fundamental limits are clear:
\begin{enumerate}
  \item \emph{Soundness is unconditional}: opacity never introduces
    false equivalences (Theorem~\ref{thm:opacity-soundness}).
  \item \emph{Completeness degrades monotonically}: as opacity
    increases, provable equivalences shrink
    (Theorem~\ref{thm:completeness-gap}).
  \item \emph{The gap is unavoidable}: by Rice's theorem, no finite
    axiom set can capture all true Python equivalences.
  \item \emph{The trade-off is navigable}: the 52-axiom sweet spot
    provides enough coverage for practical equivalence checking
    without exceeding Z3's timeout budget
    (Theorem~\ref{thm:tradeoff}).
\end{enumerate}
